% questions.tex —— 2015 年·安徽卷·文科 · 题目内容片段（无答案）
% 由 tools/generate_projects.py 依来源总账自 2015/2015_数学真题_安徽文_试题.md 生成，勿手改（改源 MD 后重跑）


\gkpart{数学}

\begin{gktip}
\textbf{来源说明}：本题面依据固定源 PDF 全 2 页逐题核对（p1＝页 933：栏 1 标题与「一、选择题」1–8（含第 7 题框图）；栏 2 第 9 题（三视图）、第 10 题（函数图象）与「二、填空题」11–15、「三、解答题」节头；栏 3 第 16、17 题（含直方图）；p2＝页 934：栏 1 第 18、19 题（含三棱锥图）、栏 2 第 20 题、栏 3 第 21 题）。原卷为 \textbf{A3 横向三栏书刊合排}，阅读序＝左→中→右；\textbf{页脚仅印合订本书页页码 933／934}，\textbf{无页眉}；标题在栏 1 内排印。\textbf{全卷不印分值、不印题量、不印「考生注意」栏}——三处节头逐字仅「一、选择题」「二、填空题」「三、解答题」；总分 150 为结构推断，\textbf{卷面无凭据}，本文件不作补写。

\textbf{源制作体系判定}（实测）：源\textbf{不是扫描件，而是 LaTeX（Computer Modern 系）重排产物}——\texttt{pdfimages -list} 零命中（全页无位图）；嵌入字体含中文 \textbf{STSong／STHeiti}，拉丁与数学 \textbf{LMRoman 系与 CM 系（CMR10／CMR7／CMMI10／CMMI7／CMSY10／CMSY7／CMEX10／MSAM10／CMMIB10）}，另有 \textbf{TeX-vect10}（向量长箭头）。文本层整体可读；根号上划线、分式横线、作答横线、\textbf{「$\parallel$」（第 15 题④）}、全部箭头与向量箭杆均为矢量路径（「$\parallel$」实测为\textbf{两条平行矢量短斜线}：各 9.98 pt、水平偏移 2.10 pt，文本层完全丢失）。

\textbf{标点与字样约定（实测）}：① \textbf{题号与句末}源内为半角 \texttt{.}，按全库约定写作全角「．」；② \textbf{正文句读}源内为\textbf{半角 ASCII} 与半角圆括号，一律归一为全角「，」「．」「：」「；」「？」「（）」（序号「①②③④⑤」为全角 STSong 出字照源）；③ \textbf{选项标号}半角「(A)」照录；④ \textbf{作答空}：填空为统一横线（实测 \textbf{31.50 pt}＝3 em，写作 \texttt{\gkblank{}}），选择题作答括注写作「（　）」（源为半角括号＋固定 10.50 pt 空档）；⑤ \textbf{分数}均为显示风格（真分式），排印一律 \(\dfrac{a}{b}\)；⑥ \textbf{不等号}按 MSAM10 实测排 \(\geqslant\)（U+2A7E）；补集号 \(\complement\) 出自 MSAM10（U+2201）；⑦ \textbf{向量}：本卷\textbf{有 CMMIB10}——第 15 题的 \(\boldsymbol{a}\)、\(\boldsymbol{b}\) 为\textbf{粗斜体}，录作 \(\boldsymbol{a}\)／\(\boldsymbol{b}\)；\(\overrightarrow{AB}\) 类为 \textbf{TeX-vect10 长箭头}照录；⑧ \textbf{正体}：\(\sin\)／\(\cos\)／\(\lg\)（LMRoman10）、虚数单位 \(\mathrm{i}\)（第 1 题）；度数记号录 \(^{\circ}\)；⑨ 全卷\textbf{句末不用「。」}（源体系特征）；词间空格照页面词界补。

\textbf{存疑与登记}（源内异常一律照录；完整清单见工作档 \texttt{figure-inventory.md} 与转录稿存疑表）：① \textbf{固定源不附参考答案}（等级 D）——解析全部由本库独立推导；② \textbf{第 12 题「\(AB=\sqrt{6}\)」的根号}在文本层整枚丢失（\texttt{pdftotext} 输出 \texttt{AB = 6}）——像素与词框双证为 \(\sqrt{6}\)，\textbf{以像素为准}；③ \textbf{第 15 题④的「$\parallel$」}为两笔矢量短斜线、文本层丢失——按语义排印 \(\parallel\)；④ \textbf{第 17 题直方图 \(a\) 柱未印数值}（纵轴仅标 0.028／0.022／0.018／0.004 与符号 \(a\)），值由面积归一约束——照源不补印；⑤ \textbf{第 19 题的图}仅含 P／A／B／C 标号，题面「\(\angle BAC=60^{\circ}\)」在图中\textbf{无对应角标注}——照源不补；⑥ 第 5、10 题作答栏「(」紧贴「是」后（排满所致）照源；⑦ 第 7 题框图为 √2 近似型新拓扑（判据 \(|a-1.414|\geqslant 0.005?\)、「是」向下、「否」向右输出）逐字照源复刻；⑧ 源内冗余与微差照源不修正（Q17 柱顶与导引线差 0.014 pt、Q19 零长度重复子路径、页 1 零长度圆帽子路径、Q9 根号上线与第 11 题分式横线等独立线段路径）。
\end{gktip}

\gkpart{试题}

\gkpart{一、选择题}

\begin{question}
\gkti{1}设 \(\mathrm{i}\) 是虚数单位，则复数 \((1-\mathrm{i})(1+2\mathrm{i})=\)（　）

(A) \(3+3\mathrm{i}\)　(B) \(-1+3\mathrm{i}\)　(C) \(3+\mathrm{i}\)　(D) \(-1+\mathrm{i}\)

\end{question}
\begin{question}
\gkti{2}设全集 \(U=\{1,2,3,4,5,6\}\)，\(A=\{1,2\}\)，\(B=\{2,3,4\}\)，则 \(A\cap(\complement_U B)=\)（　）

(A) \(\{1,2,5,6\}\)　(B) \(\{1\}\)　(C) \(\{2\}\)　(D) \(\{1,2,3,4\}\)

\end{question}
\begin{question}
\gkti{3}设 \(p\)：\(x<3\)，\(q\)：\(-1<x<3\)，则 \(p\) 是 \(q\) 成立的（　）

(A) 充分必要条件　(B) 充分不必要条件

(C) 必要不充分条件　(D) 既不充分也不必要条件

\end{question}
\begin{question}
\gkti{4}下列函数中，既是偶函数又存在零点的是（　）

(A) \(y=\cos x\)　(B) \(y=\sin x\)　(C) \(y=\ln x\)　(D) \(y=x^2+1\)

\end{question}
\begin{question}
\gkti{5}已知 \(x\)，\(y\) 满足约束条件 \(\begin{cases}x-y\geqslant 0,\\x+y-4\leqslant 0,\\y\geqslant 1,\end{cases}\) 则 \(z=-2x+y\) 的最大值是（　）

(A) \(-1\)　(B) \(-2\)　(C) \(-5\)　(D) 1

\end{question}
\begin{question}
\gkti{6}下列双曲线中，渐近线方程为 \(y=\pm 2x\) 的是（　）

(A) \(x^2-\dfrac{y^2}{4}=1\)　(B) \(\dfrac{x^2}{4}-y^2=1\)

(C) \(x^2-\dfrac{y^2}{2}=1\)　(D) \(\dfrac{x^2}{2}-y^2=1\)

\Needspace{0.4\textheight}
\end{question}
\begin{question}
\gkti{7}执行如图所示的程序框图（算法流程图），输出的 \(n\) 为（　）
\nopagebreak[4]
\begin{center}% ============================================================
% f2015ahw_q07_flowchart.tex
% 2015 安徽卷（文科）第 7 题 · 程序框图 —— 清单 F01 / §1
%
% 来源：.work/phase3-b12/figure-inventory.md §1（源 PDF 内容流逐 path 实测真页 pt）
%       源件 .work/phase3-b12/source/2015安徽文.pdf p1
%
% 坐标换算：画布 [x=1pt, y=1pt]（单位 TeX pt）。局部坐标 = 真页 pt − 原点
%   O = (125.387, 120.228)（清单 §1 几何 bbox 左下角）。真页 pt → 局部 pt 为
%   1:1 平移、不缩放——清单 §0.3 的 0.999622 是 pdftocairo 渲染件的 SVG 观测口径换算，
%   只用于与渲染件对读，不入绘图坐标。渲染层面 1 TeX pt = 0.996264 PDF pt（72/72.27），
%   整图统一 0.37%，属排版口径而非坐标换算误差。
%
% 线宽（清单 §0.4 本卷实测分档）：图形轮廓与连线 0.398pt；框内分式横线 0.42pt
% 虚线：无（本图全部实线）
% 箭头：arrows.meta Triangle 长 3.985pt、宽 2.988pt（清单 §6.3 全卷 31 枚同规格）；
%   箭尖落在清单尖端坐标；线端由 TikZ 按箭长回缩（源内线端再多深入 0.398pt = 一个
%   线宽，被填充三角完全覆盖，不可见）。
% 标签：字号 absolute 10.5pt（汉字 STSong 10.5（本图以 ctex 中文字体呈现）、数学 CMMI/CMR 10.5）；一律 anchor=base 落在清单实测基线上——
%   清单给的是字框中心：汉字字框 = 全角 em 方框（基线 = 中心 − 0.36em），拉丁/数学字框
%   = 字体 ascent/descent 方框（基线 = 中心 − 0.25em）；两式均与清单另给的基线自洽。
%
% 照源复刻的存疑项：
%    - 开始框高 16.838 / 结束框高 17.016，同宽 27.999 而高差 0.178pt —— 源内不一致，照源（B-S2）
%    - 回注折线 x 在 296.483/296.484 间有位级抖动、#22 段起讫 x 差 0.001pt（三位小数舍入），照源保留
%    - 箭头尖端与目标框之间一致留 0.199–0.282pt 间隙（源画法，非贴边），照源
%    - 分式标签 a=1+1/(1+a) 按源 600 dpi 墨迹实测逐符号落位（清单只给词块中心；尾随
%      二元算符在 TeX 数学列表中会失掉前导 glue，整块单节点会比源窄 2.3pt，故拆为逐
%      符号节点）。菱形条件按清单 §1.2 的 5 个词块中心逐块落位，块内间距由 TeX 自带；
%      本图字形（Latin Modern）与源（CM）度量略有出入，该行整体较源窄约 2.1pt，
%      块中心残差 ≤0.4pt。
% ============================================================
\begin{tikzpicture}[x=1pt, y=1pt,
  every node/.style={inner sep=0, font=\fontsize{10.5pt}{12.6pt}\selectfont},
  line cap=butt, line join=miter,
  arr/.style={-{Triangle[length=3.985pt, width=2.988pt]}},
  flow/.style={line width=0.398pt}]

  % ---- 节点（清单 §1.1 表）----
  \draw[flow, rounded corners=3.985pt] (71.041,166.194) rectangle (99.040,183.032);
  \draw[flow] (53.156,135.830) rectangle (116.924,151.623);
  \draw[flow] (158.933,102.764) -- (85.040,121.175) -- (11.146,102.764) -- (85.040,84.353) -- cycle;
  \draw[flow] (49.711,40.751) rectangle (120.369,69.698);
  % 框内分式横线（线宽 0.42pt）
  \draw[line width=0.42pt] (92.040,54.953) -- (115.673,54.953);
  \draw[flow] (59.196,11.538) rectangle (110.884,26.179);
  \draw[flow] (146.059,46.795) -- (152.195,63.654) -- (196.129,63.654) -- (189.993,46.795) -- cycle;
  \draw[flow, rounded corners=3.985pt] (157.098,15.206) rectangle (185.097,32.222);

  % ---- 连线与箭头（清单 §1.1；箭尖 = 清单尖端坐标）----
  \draw[flow, arr] (85.040,165.995) -- (85.040,151.821);
  \draw[flow, arr] (85.040,135.630) -- (85.040,121.457);
  \draw[flow, arr] (85.040,84.071) -- (85.040,69.897);
  \draw[flow, arr] (85.040,40.552) -- (85.040,26.378);
  \draw[flow, arr] (171.097,46.595) -- (171.096,32.421);
  % 「否」支路：菱形右顶点 → 右行 → 下行至平行四边形
  \draw[flow, arr] (159.215,102.764) -- (171.096,102.764) -- (171.096,63.855);
  % 回注环流：n=n+1 ↓ → 左 → 上 → 右，箭尖落在菱形上方竖直连线上
  \draw[flow, arr] (85.040,11.339) -- (85.040,0.000) -- (0.000,0.000) -- (0.000,128.544) -- (85.040,128.544);

  % ---- 标签（清单 §1.2；anchor=base 于实测基线）----
  \node[anchor=base] at (85.043,171.412) {开始};
  % a=1, n=1：按清单 §1.2 的 6 个词块中心逐块落位（同一基线 y=261.58）
  \node[anchor=base] at (59.433,141.352) {$a$};
  \node[anchor=base] at (69.203,141.352) {$=$};
  \node[anchor=base] at (80.283,141.352) {$1,$};
  \node[anchor=base] at (91.023,141.352) {$n$};
  \node[anchor=base] at (101.173,141.352) {$=$};
  \node[anchor=base] at (110.793,141.352) {$1$};
  % 菱形条件：按清单 §1.2 的 5 个词块中心逐块落位（同一基线 y=220.37）
  \node[anchor=base] at (43.773,100.142) {$|a$};
  \node[anchor=base] at (54.423,100.142) {$-$};
  \node[anchor=base] at (74.253,100.142) {$1.414|$};
  \node[anchor=base] at (94.673,100.142) {$\geqslant$};
  \node[anchor=base] at (116.103,100.142) {$0.005?$};
  \node[anchor=base] at (159.213,107.952) {否};
  \node[anchor=base] at (93.993,73.892) {是};
  % 分式标签 a=1+1/(1+a)：按源 600 dpi 墨迹实测逐符号落位（尾随二元算符
  %   在 TeX 中失 glue，整块单节点会少 2.3pt 空隙；故主行与分母均逐符号定位）
  \node[anchor=base] at (55.953,52.332) {$a$};
  \node[anchor=base] at (65.673,52.332) {$=$};
  \node[anchor=base] at (75.333,52.332) {$1$};
  \node[anchor=base] at (84.333,52.332) {$+$};
  \node[anchor=base] at (103.853,59.432) {$1$};
  \node[anchor=base] at (94.653,45.122) {$1$};
  \node[anchor=base] at (103.593,45.122) {$+$};
  \node[anchor=base] at (112.903,45.122) {$a$};
  \node[anchor=base] at (171.093,52.062) {输出$n$};
  \node[anchor=base] at (85.038,15.912) {$n=n+1$};
  \node[anchor=base] at (171.093,20.542) {结束};
\end{tikzpicture}
\end{center}
\nopagebreak[4]
(A) 3　(B) 4　(C) 5　(D) 6

\end{question}
\begin{question}
\gkti{8}直线 \(3x+4y=b\) 与圆 \(x^2+y^2-2x-2y+1=0\) 相切，则 \(b\) 的值是（　）

(A) \(-2\) 或 \(12\)　(B) \(2\) 或 \(-12\)　(C) \(-2\) 或 \(-12\)　(D) \(2\) 或 \(12\)

\Needspace{0.4\textheight}
\end{question}
\begin{question}
\gkti{9}一个四面体的三视图如图所示，则该四面体的表面积是（　）
\nopagebreak[4]
\begin{center}% ============================================================
% f2015ahw_q09_views.tex
% 2015 安徽卷（文科）第 9 题 · 三视图 —— 清单 F02 / §2
%
% 来源：.work/phase3-b12/figure-inventory.md §2（源 PDF 内容流逐 path 实测真页 pt）
%       源件 .work/phase3-b12/source/2015安徽文.pdf p1
%
% 坐标换算：画布 [x=1pt, y=1pt]（单位 TeX pt）。局部坐标 = 真页 pt − 原点
%   O = (498.358, 577.549)（清单 §2 几何 bbox 左下角）。真页 pt → 局部 pt 为
%   1:1 平移、不缩放——清单 §0.3 的 0.999622 是 pdftocairo 渲染件的 SVG 观测口径换算，
%   只用于与渲染件对读，不入绘图坐标。渲染层面 1 TeX pt = 0.996264 PDF pt（72/72.27），
%   整图统一 0.37%，属排版口径而非坐标换算误差。
%
% 线宽（清单 §0.4 本卷实测分档）：图形轮廓、尺寸线、尺寸界线 0.398pt
% 虚线：无（本图全部实线）
% 箭头：arrows.meta Triangle 长 3.985pt、宽 2.988pt（清单 §6.3 全卷 31 枚同规格）；
%   箭尖落在清单尖端坐标；线端由 TikZ 按箭长回缩（源内线端再多深入 0.398pt = 一个
%   线宽，被填充三角完全覆盖，不可见）。
% 标签：字号 absolute 10.5pt（汉字 STSong 10.5、数字/√2 CMR/CMSY 10.5）；一律 anchor=base 落在清单实测基线上——
%   清单给的是字框中心：汉字字框 = 全角 em 方框（基线 = 中心 − 0.36em），拉丁/数学字框
%   = 字体 ascent/descent 方框（基线 = 中心 − 0.25em）；两式均与清单另给的基线自洽。
%
% 照源复刻的存疑项：
%    - 俯视图右斜边的两条尺寸界线在源中为逐字节相同的两份路径（#77≡#91、#78≡#92，B-S4），
%      落墨完全重合；本图只画一次。
%    - 四种尺寸体系混用（「1」尺寸的尺寸线就在端线连线上、端线长 8.731pt；「√2」尺寸带
%      8.026pt 垂距偏移、界线长 9.878pt 且穿过尺寸线两侧各伸 4.94pt，B-S9），照源。
%    - √2 的上线（源中为独立 rule #93–#96，长 5.250pt）本图并入 $\sqrt{2}$ 字形自带的横线，
%      不另画 4 条 rule；标签位置按清单字框定基线。
% ============================================================
\begin{tikzpicture}[x=1pt, y=1pt,
  every node/.style={inner sep=0, font=\fontsize{10.5pt}{12.6pt}\selectfont},
  line cap=butt, line join=miter, line width=0.398pt,
  arr/.style={-{Triangle[length=3.985pt, width=2.988pt]}}]

  % ---- 主图形（清单 §2.1）----
  \draw (9.168,135.551) -- (96.476,135.551) -- (52.822,179.205) -- cycle;
  \draw (52.822,135.551) -- (52.822,179.205);
  \draw (140.130,135.551) -- (183.784,135.551) -- (140.130,179.205) -- cycle;
  \draw (9.168,70.070) -- (96.476,70.070) -- (52.822,26.416) -- cycle;
  \draw (52.822,26.416) -- (52.822,70.070);

  % ---- 「1」尺寸（清单 §2.2：端线 + 两段尺寸线 + 两端外向箭头）----
  % 正视图下边 ×2
  \draw (9.168,122.454) -- (9.168,131.185);
  \draw (52.822,122.454) -- (52.822,131.185);
  \draw (96.476,122.454) -- (96.476,131.185);
  \draw (12.754,126.820) -- (22.264,126.820);
  \draw (39.725,126.820) -- (49.235,126.820);
  \draw (56.408,126.820) -- (65.918,126.820);
  \draw (83.379,126.820) -- (92.889,126.820);
  \draw[arr] (12.754,126.820) -- (9.168,126.820);
  \draw[arr] (49.235,126.820) -- (52.822,126.820);
  \draw[arr] (56.408,126.820) -- (52.822,126.820);
  \draw[arr] (92.889,126.820) -- (96.476,126.820);
  % 俯视图上边 ×2
  \draw (9.168,74.436) -- (9.168,83.167);
  \draw (52.822,74.436) -- (52.822,83.167);
  \draw (96.476,74.436) -- (96.476,83.167);
  \draw (12.754,78.801) -- (22.264,78.801);
  \draw (39.725,78.801) -- (49.235,78.801);
  \draw (56.408,78.801) -- (65.918,78.801);
  \draw (83.379,78.801) -- (92.889,78.801);
  \draw[arr] (12.754,78.801) -- (9.168,78.801);
  \draw[arr] (49.235,78.801) -- (52.822,78.801);
  \draw[arr] (56.408,78.801) -- (52.822,78.801);
  \draw[arr] (92.889,78.801) -- (96.476,78.801);
  % 侧视图左边 ×1（竖直）
  \draw (127.033,135.551) -- (135.764,135.551);
  \draw (127.033,179.205) -- (135.764,179.205);
  \draw (131.399,139.137) -- (131.399,148.647);
  \draw (131.399,166.108) -- (131.399,175.618);
  \draw[arr] (131.399,139.137) -- (131.399,135.551);
  \draw[arr] (131.399,175.618) -- (131.399,179.205);
  % 侧视图下边 ×1（水平）
  \draw (140.130,122.454) -- (140.130,131.185);
  \draw (183.783,122.454) -- (183.783,131.185);
  \draw (143.716,126.820) -- (153.226,126.820);
  \draw (170.687,126.820) -- (180.197,126.820);
  \draw[arr] (143.716,126.820) -- (140.130,126.820);
  \draw[arr] (180.197,126.820) -- (183.783,126.820);

  % ---- 「√2」尺寸（清单 §2.3：界线 + 两段共线尺寸线 + 两端外向箭头，斜置）----
  % 正视图左斜边
  \draw (6.985,137.734) -- (0.000,144.719);
  \draw (50.639,181.387) -- (43.654,188.372);
  \draw (6.029,143.762) -- (16.589,154.322);
  \draw (34.050,171.783) -- (44.610,182.343);
  \draw[arr] (6.029,143.762) -- (3.493,141.226);
  \draw[arr] (44.610,182.343) -- (47.147,184.880);
  % 正视图右斜边
  \draw (98.658,137.734) -- (105.643,144.719);
  \draw (55.005,181.387) -- (61.990,188.372);
  \draw (89.055,154.322) -- (99.615,143.762);
  \draw (71.593,171.783) -- (61.033,182.343);
  \draw[arr] (99.615,143.762) -- (102.151,141.226);
  \draw[arr] (61.033,182.343) -- (58.497,184.880);
  % 俯视图左斜边
  \draw (6.985,67.887) -- (0.000,60.902);
  \draw (50.639,24.233) -- (43.654,17.248);
  \draw (16.589,51.298) -- (6.029,61.858);
  \draw (34.050,33.837) -- (44.610,23.277);
  \draw[arr] (6.029,61.858) -- (3.493,64.395);
  \draw[arr] (44.610,23.277) -- (47.147,20.741);
  % 俯视图右斜边
  \draw (98.658,67.887) -- (105.643,60.902);
  \draw (55.005,24.233) -- (61.990,17.248);
  \draw (89.055,51.298) -- (99.615,61.858);
  \draw (71.593,33.837) -- (61.033,23.277);
  \draw[arr] (99.615,61.858) -- (102.151,64.395);
  \draw[arr] (61.033,23.277) -- (58.497,20.741);

  % ---- 标签（清单 §2.4；anchor=base 于实测基线）----
  \node[anchor=base] at (52.822,102.301) {正(主)视图};
  \node[anchor=base] at (161.952,102.271) {侧(左)视图};
  \node[anchor=base] at (52.822,1.471) {俯视图};
  \node[anchor=base] at (30.992,123.436) {$1$};
  \node[anchor=base] at (74.652,123.436) {$1$};
  \node[anchor=base] at (30.992,75.416) {$1$};
  \node[anchor=base] at (74.652,75.416) {$1$};
  \node[anchor=base] at (131.402,153.996) {$1$};
  \node[anchor=base] at (161.952,123.436) {$1$};
  \node[anchor=base] at (20.077,159.856) {$\sqrt{2}$};
  \node[anchor=base] at (79.017,159.856) {$\sqrt{2}$};
  \node[anchor=base] at (20.077,37.626) {$\sqrt{2}$};
  \node[anchor=base] at (79.017,37.626) {$\sqrt{2}$};
\end{tikzpicture}
\end{center}
\nopagebreak[4]
(A) \(1+\sqrt{3}\)　(B) \(2+\sqrt{3}\)　(C) \(1+2\sqrt{2}\)　(D) \(2\sqrt{2}\)

\Needspace{0.4\textheight}
\end{question}
\begin{question}
\gkti{10}函数 \(f(x)=ax^3+bx^2+cx+d\) 的图象如图所示，则下列结论成立的是（　）
\nopagebreak[4]
\begin{center}% ============================================================
% f2015ahw_q10_graph.tex
% 2015 安徽卷（文科）第 10 题 · 三次函数图象 —— 清单 F03 / §3
%
% 来源：.work/phase3-b12/figure-inventory.md §3（源 PDF 内容流逐 path 实测真页 pt）
%       源件 .work/phase3-b12/source/2015安徽文.pdf p1
%
% 坐标换算：画布 [x=1pt, y=1pt]（单位 TeX pt）。局部坐标 = 真页 pt − 原点
%   O = (507.935, 380.106)（清单 §3 几何 bbox 左下角）。真页 pt → 局部 pt 为
%   1:1 平移、不缩放——清单 §0.3 的 0.999622 是 pdftocairo 渲染件的 SVG 观测口径换算，
%   只用于与渲染件对读，不入绘图坐标。渲染层面 1 TeX pt = 0.996264 PDF pt（72/72.27），
%   整图统一 0.37%，属排版口径而非坐标换算误差。
%
% 线宽（清单 §0.4 本卷实测分档）：坐标轴与虚线 0.398pt；函数曲线 0.797pt（全卷唯一加粗曲线）
% 虚线：[2.989 2.989] 相位 0（清单 §0.4 本卷唯一虚线规格）
% 箭头：arrows.meta Triangle 长 3.985pt、宽 2.988pt（清单 §6.3 全卷 31 枚同规格）；
%   箭尖落在清单尖端坐标；线端由 TikZ 按箭长回缩（源内线端再多深入 0.398pt = 一个
%   线宽，被填充三角完全覆盖，不可见）。
% 标签：字号 absolute 10.5pt（数学斜体 CMMI 10.5（下标 2/1 为 CMR7 6.974））；一律 anchor=base 落在清单实测基线上——
%   清单给的是字框中心：汉字字框 = 全角 em 方框（基线 = 中心 − 0.36em），拉丁/数学字框
%   = 字体 ascent/descent 方框（基线 = 中心 − 0.25em）；两式均与清单另给的基线自洽。
%
% 照源复刻的存疑项：
%    - 函数曲线在源中为 299 段三次 Bézier 的连续链（1×v + 297×c + 1×y，无 l、无 h），
%      逐段对弦垂距 ≤0.0023pt（线宽 0.797pt 的 0.29%）——几何上是折线。本图按折线复刻：
%      300 个链节点连线（清单 §6.4「还原时应按折线（300 个节点）处理」）；节点坐标取源内容流
%      实测真页 pt（四位小数，与清单三位小数的舍入差 ≤0.001pt）。
%    - 两条虚线均未触到曲线：x₁ 处顶 (554.707,469.161) 距曲线 1.691pt、x₂ 处底 (620.189,402.707)
%      距曲线 1.197pt，两处空隙量不等（B-S11），照源不接；虚线绘制方向照源（#106 自上而下）。
%    - P 点无标记点（全图仅两枚箭头为填充），照源不补圆点。
% ============================================================
\begin{tikzpicture}[x=1pt, y=1pt,
  every node/.style={inner sep=0, font=\fontsize{10.5pt}{12.6pt}\selectfont},
  line cap=butt, line join=miter,
  arr/.style={-{Triangle[length=3.985pt, width=2.988pt]}},
  dashline/.style={line width=0.398pt, dash pattern=on 2.989pt off 2.989pt}]

  % ---- 坐标轴（x 轴 164.794pt、y 轴 127.376pt）与箭头 ----
  \draw[line width=0.398pt, arr] (0.000,56.127) -- (168.381,56.127);
  \draw[line width=0.398pt, arr] (37.418,0.000) -- (37.418,130.963);

  % ---- 两条竖直虚线（x₂ 自曲线极小点上行、x₁ 自 x 轴上行）----
  \draw[dashline] (112.254,22.601) -- (112.254,56.127);
  \draw[dashline] (46.772,89.055) -- (46.772,56.127);

  % ---- 三次函数曲线 y=f(x)：源 299 段 Bézier 链的 300 个节点，折线复刻（线宽 0.797pt）----
  \draw[line width=0.797pt]
    (13.0966,15.0395) -- (13.5406,17.2795) -- (13.9846,19.4805) -- (14.4296,21.6415) --
    (14.8736,23.7635) -- (15.3176,25.8455) -- (15.7616,27.8895) -- (16.2056,29.8945) --
    (16.6506,31.8625) -- (17.0946,33.7905) -- (17.5386,35.6815) -- (17.9826,37.5365) --
    (18.4266,39.3535) -- (18.8716,41.1335) -- (19.3156,42.8765) -- (19.7596,44.5845) --
    (20.2036,46.2555) -- (20.6476,47.8905) -- (21.0926,49.4905) -- (21.5366,51.0545) --
    (21.9806,52.5845) -- (22.4246,54.0775) -- (22.8686,55.5375) -- (23.3136,56.9635) --
    (23.7576,58.3575) -- (24.2016,59.7145) -- (24.6456,61.0405) -- (25.0896,62.3335) --
    (25.5346,63.5915) -- (25.9786,64.8195) -- (26.4226,66.0135) -- (26.8666,67.1745) --
    (27.3106,68.3065) -- (27.7556,69.4055) -- (28.1996,70.4745) -- (28.6436,71.5115) --
    (29.0876,72.5175) -- (29.5316,73.4945) -- (29.9766,74.4395) -- (30.4206,75.3545) --
    (30.8646,76.2425) -- (31.3086,77.1005) -- (31.7526,77.9275) -- (32.1976,78.7285) --
    (32.6416,79.4995) -- (33.0856,80.2415) -- (33.5296,80.9565) -- (33.9736,81.6435) --
    (34.4186,82.3035) -- (34.8626,82.9365) -- (35.3066,83.5425) -- (35.7506,84.1205) --
    (36.1946,84.6735) -- (36.6396,85.2005) -- (37.0836,85.7015) -- (37.5276,86.1765) --
    (37.9716,86.6265) -- (38.4156,87.0525) -- (38.8606,87.4525) -- (39.3046,87.8295) --
    (39.7486,88.1805) -- (40.1926,88.5085) -- (40.6366,88.8125) -- (41.0816,89.0935) --
    (41.5256,89.3525) -- (41.9696,89.5875) -- (42.4136,89.8005) -- (42.8576,89.9905) --
    (43.3026,90.1585) -- (43.7466,90.3055) -- (44.1906,90.4305) -- (44.6346,90.5355) --
    (45.0786,90.6175) -- (45.5236,90.6795) -- (45.9676,90.7225) -- (46.4116,90.7435) --
    (46.8556,90.7465) -- (47.2996,90.7295) -- (47.7446,90.6925) -- (48.1886,90.6375) --
    (48.6326,90.5625) -- (49.0766,90.4695) -- (49.5206,90.3585) -- (49.9656,90.2285) --
    (50.4096,90.0815) -- (50.8536,89.9185) -- (51.2976,89.7365) -- (51.7416,89.5385) --
    (52.1866,89.3245) -- (52.6306,89.0925) -- (53.0746,88.8455) -- (53.5186,88.5825) --
    (53.9626,88.3045) -- (54.4076,88.0105) -- (54.8516,87.7015) -- (55.2956,87.3785) --
    (55.7396,87.0405) -- (56.1836,86.6885) -- (56.6286,86.3225) -- (57.0726,85.9435) --
    (57.5166,85.5505) -- (57.9606,85.1445) -- (58.4056,84.7265) -- (58.8496,84.2955) --
    (59.2936,83.8515) -- (59.7376,83.3955) -- (60.1816,82.9285) -- (60.6266,82.4505) --
    (61.0706,81.9605) -- (61.5146,81.4595) -- (61.9586,80.9485) -- (62.4026,80.4265) --
    (62.8476,79.8945) -- (63.2916,79.3535) -- (63.7356,78.8015) -- (64.1796,78.2415) --
    (64.6236,77.6715) -- (65.0686,77.0945) -- (65.5126,76.5055) -- (65.9566,75.9115) --
    (66.4006,75.3085) -- (66.8446,74.6975) -- (67.2896,74.0805) -- (67.7336,73.4545) --
    (68.1776,72.8235) -- (68.6216,72.1855) -- (69.0656,71.5395) -- (69.5106,70.8885) --
    (69.9546,70.2315) -- (70.3986,69.5705) -- (70.8426,68.9025) -- (71.2866,68.2305) --
    (71.7316,67.5525) -- (72.1756,66.8715) -- (72.6196,66.1865) -- (73.0636,65.4955) --
    (73.5076,64.8035) -- (73.9526,64.1065) -- (74.3966,63.4065) -- (74.8406,62.7045) --
    (75.2846,61.9995) -- (75.7286,61.2925) -- (76.1736,60.5845) -- (76.6176,59.8725) --
    (77.0616,59.1615) -- (77.5056,58.4475) -- (77.9496,57.7335) -- (78.3946,57.0175) --
    (78.8386,56.3025) -- (79.2826,55.5875) -- (79.7266,54.8725) -- (80.1706,54.1575) --
    (80.6156,53.4435) -- (81.0596,52.7315) -- (81.5036,52.0195) -- (81.9476,51.3085) --
    (82.3916,50.5995) -- (82.8366,49.8935) -- (83.2806,49.1885) -- (83.7246,48.4875) --
    (84.1686,47.7885) -- (84.6126,47.0925) -- (85.0576,46.4005) -- (85.5016,45.7125) --
    (85.9456,45.0275) -- (86.3896,44.3485) -- (86.8336,43.6715) -- (87.2786,43.0015) --
    (87.7226,42.3365) -- (88.1666,41.6755) -- (88.6106,41.0205) -- (89.0546,40.3715) --
    (89.4996,39.7285) -- (89.9436,39.0915) -- (90.3876,38.4615) -- (90.8316,37.8385) --
    (91.2756,37.2235) -- (91.7206,36.6145) -- (92.1646,36.0135) -- (92.6086,35.4215) --
    (93.0526,34.8375) -- (93.4966,34.2615) -- (93.9416,33.6945) -- (94.3856,33.1365) --
    (94.8296,32.5885) -- (95.2736,32.0495) -- (95.7176,31.5205) -- (96.1626,31.0025) --
    (96.6066,30.4935) -- (97.0506,29.9965) -- (97.4946,29.5095) -- (97.9386,29.0345) --
    (98.3836,28.5705) -- (98.8276,28.1185) -- (99.2716,27.6795) -- (99.7156,27.2515) --
    (100.1596,26.8375) -- (100.6046,26.4355) -- (101.0486,26.0465) -- (101.4926,25.6705) --
    (101.9366,25.3085) -- (102.3806,24.9615) -- (102.8256,24.6285) -- (103.2696,24.3085) --
    (103.7136,24.0045) -- (104.1576,23.7155) -- (104.6016,23.4425) -- (105.0466,23.1835) --
    (105.4906,22.9405) -- (105.9346,22.7145) -- (106.3786,22.5055) -- (106.8236,22.3125) --
    (107.2676,22.1355) -- (107.7116,21.9775) -- (108.1556,21.8355) -- (108.5996,21.7105) --
    (109.0436,21.6045) -- (109.4886,21.5175) -- (109.9326,21.4495) -- (110.3766,21.3985) --
    (110.8206,21.3675) -- (111.2656,21.3565) -- (111.7096,21.3645) -- (112.1536,21.3925) --
    (112.5976,21.4415) -- (113.0416,21.5085) -- (113.4866,21.5985) -- (113.9306,21.7095) --
    (114.3746,21.8405) -- (114.8186,21.9945) -- (115.2626,22.1695) -- (115.7076,22.3655) --
    (116.1516,22.5855) -- (116.5956,22.8275) -- (117.0396,23.0925) -- (117.4836,23.3815) --
    (117.9286,23.6915) -- (118.3726,24.0275) -- (118.8166,24.3865) -- (119.2606,24.7695) --
    (119.7046,25.1785) -- (120.1496,25.6115) -- (120.5936,26.0695) -- (121.0376,26.5515) --
    (121.4816,27.0605) -- (121.9256,27.5955) -- (122.3706,28.1555) -- (122.8146,28.7425) --
    (123.2586,29.3565) -- (123.7026,29.9975) -- (124.1466,30.6655) -- (124.5916,31.3615) --
    (125.0356,32.0845) -- (125.4796,32.8355) -- (125.9236,33.6155) -- (126.3676,34.4235) --
    (126.8126,35.2605) -- (127.2566,36.1265) -- (127.7006,37.0215) -- (128.1446,37.9475) --
    (128.5886,38.9025) -- (129.0336,39.8865) -- (129.4776,40.9035) -- (129.9216,41.9495) --
    (130.3656,43.0275) -- (130.8096,44.1355) -- (131.2546,45.2765) -- (131.6986,46.4485) --
    (132.1426,47.6515) -- (132.5866,48.8895) -- (133.0306,50.1575) -- (133.4756,51.4605) --
    (133.9196,52.7945) -- (134.3636,54.1635) -- (134.8076,55.5655) -- (135.2516,57.0015) --
    (135.6966,58.4725) -- (136.1406,59.9775) -- (136.5846,61.5175) -- (137.0286,63.0925) --
    (137.4726,64.7025) -- (137.9176,66.3495) -- (138.3616,68.0305) -- (138.8056,69.7485) --
    (139.2496,71.5035) -- (139.6936,73.2945) -- (140.1386,75.1225) -- (140.5826,76.9875) --
    (141.0266,78.8905) -- (141.4706,80.8305) -- (141.9146,82.8085) -- (142.3596,84.8255) --
    (142.8036,86.8815) -- (143.2476,88.9755) -- (143.6916,91.1085) -- (144.1356,93.2815) --
    (144.5806,95.4945) -- (145.0246,97.7465) -- (145.4686,100.0385) -- (145.9126,102.3725);

  % ---- 标签（清单 §3.2；anchor=base，下标随主字同点）----
  \node[anchor=base] at (30.765,129.719) {$y$};
  \node[anchor=base] at (28.895,94.449) {$P$};
  \node[anchor=base] at (112.005,61.399) {$x_2$};
  \node[anchor=base] at (46.520,47.909) {$x_1$};
  \node[anchor=base] at (29.425,45.249) {$O$};
  \node[anchor=base] at (168.385,47.909) {$x$};
\end{tikzpicture}
\end{center}
\nopagebreak[4]
(A) \(a>0\)，\(b<0\)，\(c>0\)，\(d>0\)　(B) \(a>0\)，\(b<0\)，\(c<0\)，\(d>0\)

(C) \(a<0\)，\(b<0\)，\(c>0\)，\(d>0\)　(D) \(a>0\)，\(b>0\)，\(c>0\)，\(d<0\)

\end{question}
\gkpart{二、填空题}

\begin{question}
\gkti{11}计算：\(\lg\dfrac{5}{2}+2\lg 2-\left(\dfrac{1}{2}\right)^{-1}=\)\gkblank{}．

\end{question}
\begin{question}
\gkti{12}在 \(\triangle ABC\) 中，\(AB=\sqrt{6}\)，\(\angle A=75^{\circ}\)，\(\angle B=45^{\circ}\)，则 \(AC=\)\gkblank{}．

\end{question}
\begin{question}
\gkti{13}已知数列 \(\{a_n\}\) 中，\(a_1=1\)，\(a_n=a_{n-1}+\dfrac{1}{2}\)（\(n\geqslant 2\)），则数列 \(\{a_n\}\) 的前 9 项和等于\gkblank{}．

\end{question}
\begin{question}
\gkti{14}在平面直角坐标系 \(xOy\) 中，若直线 \(y=2a\) 与函数 \(y=|x-a|-1\) 的图象只有一个交点，则 \(a\) 的值为\gkblank{}．

\end{question}
\begin{question}
\gkti{15}\(\triangle ABC\) 是边长为 2 的等边三角形，已知向量 \(\boldsymbol{a}\)，\(\boldsymbol{b}\) 满足 \(\overrightarrow{AB}=2\boldsymbol{a}\)，\(\overrightarrow{AC}=2\boldsymbol{a}+\boldsymbol{b}\)，则下列结论中正确的是\gkblank{}．（写出所有正确结论的序号）

① \(\boldsymbol{a}\) 为单位向量；　② \(\boldsymbol{b}\) 为单位向量；　③ \(\boldsymbol{a}\perp\boldsymbol{b}\)；

④ \(\boldsymbol{b}\parallel\overrightarrow{BC}\)；　⑤ \((4\boldsymbol{a}+\boldsymbol{b})\perp\overrightarrow{BC}\)．

\end{question}
\gkpart{三、解答题}

\begin{question}
\gkti{16}已知函数 \(f(x)=(\sin x+\cos x)^2+\cos 2x\)．

（1） 求 \(f(x)\) 最小正周期；

（2） 求 \(f(x)\) 在区间 \(\left[0,\dfrac{\pi}{2}\right]\) 上的最大值和最小值．

\Needspace{0.4\textheight}
\end{question}
\begin{question}
\gkti{17}某企业为了解下属某部门对本企业职工的服务情况，随机访问 50 名职工，根据这 50 名职工对该部门的评分，绘制频率分布直方图（如图所示），其中样本数据分组区间为：\([40,50)\)，\([50,60)\)，\(\cdots\)，\([80,90)\)，\([90,100]\)．
\nopagebreak[4]
\begin{center}% ============================================================
% f2015ahw_q17_histogram.tex
% 2015 安徽卷（文科）第 17 题 · 频率分布直方图 —— 清单 F04 / §4
%
% 来源：.work/phase3-b12/figure-inventory.md §4（源 PDF 内容流逐 path 实测真页 pt）
%       源件 .work/phase3-b12/source/2015安徽文.pdf p1
%
% 坐标换算：画布 [x=1pt, y=1pt]（单位 TeX pt）。局部坐标 = 真页 pt − 原点
%   O = (866.533, 257.207)（清单 §4 几何 bbox 左下角）。真页 pt → 局部 pt 为
%   1:1 平移、不缩放——清单 §0.3 的 0.999622 是 pdftocairo 渲染件的 SVG 观测口径换算，
%   只用于与渲染件对读，不入绘图坐标。渲染层面 1 TeX pt = 0.996264 PDF pt（72/72.27），
%   整图统一 0.37%，属排版口径而非坐标换算误差。
%
% 线宽（清单 §0.4 本卷实测分档）：坐标轴、柱体、虚线导引线 0.398pt；纵轴分式横线 0.42pt
% 虚线：[2.989 2.989] 相位 0（5 条导引线）
% 箭头：arrows.meta Triangle 长 3.985pt、宽 2.988pt（清单 §6.3 全卷 31 枚同规格）；
%   箭尖落在清单尖端坐标；线端由 TikZ 按箭长回缩（源内线端再多深入 0.398pt = 一个
%   线宽，被填充三角完全覆盖，不可见）。
% 标签：字号 absolute 10.5pt（汉字 STSong 10.5、刻度数字 CMR/LMRoman 10.5、纵轴 a 为 CMMI 10.5 斜体）；一律 anchor=base 落在清单实测基线上——
%   清单给的是字框中心：汉字字框 = 全角 em 方框（基线 = 中心 − 0.36em），拉丁/数学字框
%   = 字体 ascent/descent 方框（基线 = 中心 − 0.25em）；两式均与清单另给的基线自洽。
%
% 照源复刻的存疑项：
%    - 6 柱（不是 8 柱；8 是横轴刻度个数），相邻柱各自成矩形、共用边两处相差 0.001pt（B-S3b），
%      照源逐柱独立画，不合并成连续条带、不改柱宽。
%    - a 柱顶 295.434 与其等高虚线导引线 295.448 相差 0.014pt（B-S3）——源内不一致，照源不对齐。
%    - 0→40 的轴长被压缩为 19.843pt（1 格，正常 4 格），其间以 3 峰 3 谷 zigzag 占位（B-S10）；
%      柱 [40,50) 的左边界 913.994 与刻度「40」文字中心重合，照源。
%    - 断轴 zigzag 与 x 轴主干在源中为同一条 14 节点折线（#119），本图按折线照录。
% ============================================================
\begin{tikzpicture}[x=1pt, y=1pt,
  every node/.style={inner sep=0, font=\fontsize{10.5pt}{12.6pt}\selectfont},
  line cap=butt, line join=miter,
  arr/.style={-{Triangle[length=3.985pt, width=2.988pt]}},
  dashline/.style={line width=0.398pt, dash pattern=on 2.989pt off 2.989pt}]

  % ---- 坐标轴：单条 14 节点折线（含断轴 zigzag），y 轴 145.251pt、x 轴主干 145.233pt ----
  \draw[line width=0.398pt] (27.618,157.980) -- (27.618,12.729) -- (31.586,12.729) -- (32.579,14.004) -- (34.563,11.453) -- (35.555,12.729) -- (36.547,14.004) -- (38.532,11.453) -- (39.524,12.728) -- (40.516,14.004) -- (42.501,11.453) -- (43.493,12.728) -- (47.461,12.728) -- (192.694,12.729);
  \draw[line width=0.398pt, arr] (27.618,157.980) -- (27.618,161.567);
  \draw[line width=0.398pt, arr] (192.694,12.729) -- (196.281,12.729);
  % 纵轴「频率/组距」分式横线（线宽 0.42pt）
  \draw[line width=0.42pt] (32.512,160.887) -- (53.512,160.887);

  % ---- 6 柱（各宽 19.843pt；逐柱独立矩形，共用边照源微差）----
  \draw (47.461,12.729) rectangle (67.304,29.728);
  \draw (67.303,12.729) rectangle (87.146,38.227);
  \draw (87.146,12.729) rectangle (106.989,106.288);
  \draw (106.989,12.729) rectangle (126.832,131.786);
  \draw (126.831,12.729) rectangle (146.674,106.288);
  \draw (146.674,12.729) rectangle (166.517,89.289);

  % ---- 5 条虚线导引线（起于纵轴 x=894.151，终于各自柱左上角）----
  \draw[dashline] (27.618,29.728) -- (47.461,29.728);
  \draw[dashline] (27.618,38.241) -- (67.303,38.241);
  \draw[dashline] (27.618,106.288) -- (126.831,106.288);
  \draw[dashline] (27.618,131.786) -- (106.989,131.786);
  \draw[dashline] (27.618,89.289) -- (146.674,89.289);

  % ---- 标签（清单 §4.4/§4.5；anchor=base 于实测基线）----
  \node[anchor=base] at (27.617,2.258) {$0$};
  \node[anchor=base] at (47.457,2.258) {$40$};
  \node[anchor=base] at (67.307,2.258) {$50$};
  \node[anchor=base] at (87.147,2.258) {$60$};
  \node[anchor=base] at (106.987,2.258) {$70$};
  \node[anchor=base] at (126.827,2.258) {$80$};
  \node[anchor=base] at (146.677,2.258) {$90$};
  \node[anchor=base] at (166.517,2.258) {$100$};
  \node[anchor=base] at (43.007,165.363) {频率};
  \node[anchor=base] at (43.007,151.063) {组距};
  \node[anchor=base] at (11.957,128.618) {$0.028$};
  \node[anchor=base] at (11.957,103.128) {$0.022$};
  \node[anchor=base] at (11.957,86.128) {$0.018$};
  \node[anchor=base] at (21.147,35.978) {$a$};
  \node[anchor=base] at (11.957,26.508) {$0.004$};
  \node[anchor=base] at (196.277,18.243) {分数};
\end{tikzpicture}
\end{center}
\nopagebreak[4]
（1） 求频率分布直方图中 \(a\) 的值；

（2） 估计该企业的职工对该部门评分不低于 80 的概率；

（3） 从评分在 \([40,60)\) 的受访职工中，随机抽取 2 人，求此 2 人评分都在 \([40,50)\) 的概率．

\end{question}
\begin{question}
\gkti{18}已知数列 \(\{a_n\}\) 是递增的等比数列，且 \(a_1+a_4=9\)，\(a_2a_3=8\)．

（1） 求数列 \(\{a_n\}\) 的通项公式；

（2） 设 \(S_n\) 为数列 \(\{a_n\}\) 的前 \(n\) 项和，\(b_n=\dfrac{a_{n+1}}{S_nS_{n+1}}\)，求数列 \(\{b_n\}\) 的前 \(n\) 项和 \(T_n\)．

\Needspace{0.4\textheight}
\end{question}
\begin{question}
\gkti{19}如图，三棱锥 \(P\text{-}ABC\) 中，\(PA\perp\) 平面 \(ABC\)，\(PA=1\)，\(AB=1\)，\(AC=2\)，\(\angle BAC=60^{\circ}\)．
\nopagebreak[4]
\begin{center}% ============================================================
% f2015ahw_q19_pyramid.tex
% 2015 安徽卷（文科）第 19 题 · 三棱锥 P-ABC —— 清单 F05 / §5
%
% 来源：.work/phase3-b12/figure-inventory.md §5（源 PDF 内容流逐 path 实测真页 pt）
%       源件 .work/phase3-b12/source/2015安徽文.pdf p2
%
% 坐标换算：画布 [x=1pt, y=1pt]（单位 TeX pt）。局部坐标 = 真页 pt − 原点
%   O = (222.468, 240.178)（清单 §5 几何 bbox 左下角）。真页 pt → 局部 pt 为
%   1:1 平移、不缩放——清单 §0.3 的 0.999622 是 pdftocairo 渲染件的 SVG 观测口径换算，
%   只用于与渲染件对读，不入绘图坐标。渲染层面 1 TeX pt = 0.996264 PDF pt（72/72.27），
%   整图统一 0.37%，属排版口径而非坐标换算误差。
%
% 线宽（清单 §0.4 本卷实测分档）：6 条棱一律 0.3985pt（p2 档）
% 虚线：[2.98883 2.98883] 相位 0（仅 AC 棱）
% 箭头：无（本图不含箭头）
% 标签：字号 absolute 10.5pt（顶点字母 CMMI 10.5 数学斜体）；一律 anchor=base 落在清单实测基线上——
%   清单给的是字框中心：汉字字框 = 全角 em 方框（基线 = 中心 − 0.36em），拉丁/数学字框
%   = 字体 ascent/descent 方框（基线 = 中心 − 0.25em）；两式均与清单另给的基线自洽。
%
% 照源复刻的存疑项：
%    - AC 为唯一虚线棱（源 [2.98883 2.98883] 0，与 p1 的 2.989 为同一量的不同舍入，B-S7）；
%      其余五棱实线，线宽 0.3985pt（p2 独有档位，比 p1 的 0.398 多 0.0005pt），照源不复核。
%    - 题面含 ∠BAC = 60°，但源图内无任何角度标注（无角弧、无度数文字、无直角符号、
%      无 BM 线段与 M 点，B-S6）——照源不补标注。
%    - 源每条棱的路径都含 1 个零长度重复子路径（x y m x y m … l S，B-S5），不落墨，已丢弃。
% ============================================================
\begin{tikzpicture}[x=1pt, y=1pt,
  every node/.style={inner sep=0, font=\fontsize{10.5pt}{12.6pt}\selectfont},
  line cap=butt, line join=miter, line width=0.3985pt,
  hidden/.style={dash pattern=on 2.98883pt off 2.98883pt}]

  % ---- 棱（清单 §5.1；PA 竖直、BC 水平）：PA / PB / PC / AB / BC 实线，AC 虚线 ----
  \draw (10.911,113.385) -- (10.911,37.412);
  \draw (10.911,113.385) -- (37.767,10.556);
  \draw (10.911,113.385) -- (169.357,10.556);
  \draw (10.911,37.412) -- (37.767,10.556);
  \draw (37.767,10.556) -- (169.357,10.556);
  \draw[hidden] (10.911,37.412) -- (169.357,10.556);

  % ---- 标签（清单 §5.2；anchor=base 于实测基线）----
  \node[anchor=base] at (10.182,116.767) {$P$};
  \node[anchor=base] at (3.942,33.827) {$A$};
  \node[anchor=base] at (32.572,2.037) {$B$};
  \node[anchor=base] at (173.912,2.037) {$C$};
\end{tikzpicture}
\end{center}
\nopagebreak[4]
（1） 求三棱锥 \(P\text{-}ABC\) 的体积；

（2） 证明：在线段 \(PC\) 上存在点 \(M\)，使得 \(AC\perp BM\)，并求 \(\dfrac{PM}{MC}\) 的值．

\end{question}
\begin{question}
\gkti{20}设椭圆 \(E\) 的方程为 \(\dfrac{x^2}{a^2}+\dfrac{y^2}{b^2}=1\)（\(a>b>0\)），点 \(O\) 为坐标原点，点 \(A\) 的坐标为 \((a,0)\)，点 \(B\) 的坐标为 \((0,b)\)，点 \(M\) 在线段 \(AB\) 上，满足 \(|BM|=2|MA|\)，直线 \(OM\) 的斜率为 \(\dfrac{\sqrt{5}}{10}\)．

（1） 求 \(E\) 的离心率 \(e\)；

（2） 设点 \(C\) 的坐标为 \((0,-b)\)，\(N\) 为线段 \(AC\) 的中点，证明：\(MN\perp AB\)．

\end{question}
\begin{question}
\gkti{21}已知函数 \(f(x)=\dfrac{ax}{(x+r)^2}\)（\(a>0\)，\(r>0\)）．

（1） 求 \(f(x)\) 的定义域，并讨论 \(f(x)\) 的单调性；

（2） 若 \(\dfrac{a}{r}=400\)，求 \(f(x)\) 在 \((0,+\infty)\) 内的极值．

\end{question}